%%%%%%%%%%%%%%%%%%%%%%%%%%%%%%%%%
% Abstract                      %
%%%%%%%%%%%%%%%%%%%%%%%%%%%%%%%%%

\section*{Abstract}
\addcontentsline{toc}{section}{Abstract}

\vspace{1.4cm}
\large
\noindent

The BAKO Motors Intelligent Supervision System (ISS) is a full-stack embedded IoT platform designed and delivered as part of the ISS~496 senior project at MedTech Mediterranean Institute of Technology, in direct collaboration with BAKO Motors as the industrial partner. The project addresses five critical operational gaps in BAKO Motors' electric vehicle prototype: the absence of real-time sensor coverage beyond the battery pack, the restriction of CAN bus communication to the BMS only, the lack of a monitoring interface for engineers, insufficient system documentation for maintenance teams, and the inability to perform remote diagnostics.

The system is built around a custom two-layer PCB integrating an ESP32 microcontroller, an MCP2515/TJA1051T CAN controller and transceiver, three DS18B20 temperature sensors, two ACS712 current sensors, and a SIM800L GSM/GPRS module. The firmware implements a single-threaded cooperative loop that decodes SAE~J1939 CAN frames from the O'CELL IFS60.8-500-F-E3 BMS, polls external sensors at appropriate rates, and publishes structured telemetry snapshots over both Wi-Fi and cellular. A Python-based seven-scenario sensor emulator was developed to enable full system validation without physical hardware. A real-time web dashboard streams live vehicle telemetry at 10~Hz via WebSocket and supports XLSX data export. A cloud pipeline transmits telemetry over GPRS to a FastAPI server and SQLite database hosted on a Linux VPS, with live browser access via a dedicated WebSocket endpoint.

All defined V\&V acceptance criteria were met: CAN frame reception exceeded 99\% after termination correction, DS18B20 temperature accuracy was $\pm 0.5$\textdegree{}C, ACS712 current measurement error was $\pm 1.8\%$, voltage measurement error was $\pm 0.9\%$, GSM alert latency was under 15~seconds, and zero cloud snapshots were lost over a one-hour sustained test. The PCB was fabricated, fully populated, installed in the vehicle chassis, and field-validated. The project was completed within an eight-sprint Agile cycle at a total cost of 240~TND, representing a 5.9\% underspend against the planned budget.\\

\textbf{Keywords:} Electric Vehicle, CAN Bus, SAE J1939, ESP32, BMS, LiFePO4, IoT, Real-Time Monitoring, Telemetry, GSM/GPRS, FastAPI, PCB Design, BAKO Motors.
